\documentclass[11pt,a4paper]{article}
\usepackage[top=0.75in,bottom=0.75in,left=0.75in,right=0.75in]{geometry}
\usepackage[T1]{fontenc}
\usepackage{times,microtype,enumitem,hyperref,titlesec}
\pagestyle{empty}
\setlength{\parindent}{0pt}
\setlength{\parskip}{0pt}
\linespread{1.10}
\hypersetup{hidelinks,breaklinks=true}
\titleformat{\section}{\bfseries\normalsize}{}{0pt}{}[\titlerule]
\titlespacing*{\section}{0pt}{10pt}{4pt}
\newcommand{\entry}[3]{\noindent\textbf{#1}\hfill\textbf{#3}\\[-1pt]\textit{#2}\par}
\newenvironment{cvitems}{\begin{itemize}[leftmargin=1.1em,label=\textbullet,itemsep=2pt,topsep=3pt,parsep=0pt,partopsep=0pt]}{\end{itemize}}
\begin{document}
{\LARGE\bfseries Aarush Ram Anandh}\par\vspace{4pt}
\href{mailto:aarush.anandh_ug2024@ashoka.edu.in}{aarush.anandh\_ug2024@ashoka.edu.in}\quad|\quad +91 9659786783\par\vspace{2pt}
\href{https://github.com/MushuRam01}{github.com/MushuRam01}\quad|\quad
\href{https://www.linkedin.com/in/aarush-ram-anandh/}{linkedin.com/in/aarush-ram-anandh}\par
\vspace{20pt}
11 November 2026\par\vspace{12pt}
Hiring Team\\
Honeywell\par\vspace{12pt}
\textbf{Application: AI Internship}\par\vspace{14pt}
Dear Hiring Team,\par\vspace{10pt}
I am applying for the AI internship at Honeywell. I am completing a B.Sc. (Hons.) in Computer Science and Mathematics at Ashoka University, expected in 2028. I am interested in the practical work of building AI algorithms around sensors, understanding failure cases and turning models into systems that people can inspect and use.\par\vspace{10pt}
At Avartan Labs, I raised rare-class correct detection from 40\% to 80\%, a gain of 40 percentage points, in an RT-DETR detection, classification and tracking pipeline for 18 waste-material classes. I investigated focal loss, logit adjustment, class-specific thresholds and targeted fine-tuning to address rare-class performance, and used crop-based refinement to examine detector failures. I also helped set up a near-infrared waste-segregation system for visually similar plastics, including parts procurement and preliminary experiments. This work is directly relevant to my interest in industrial AI and automation.\par\vspace{10pt}
My Shooting Stability Trainer project combines OpenCV marker tracking with a Flask-SocketIO server and a JavaScript interface. It provides segmentation-mask feedback, scale calibration and manually recorded motion traces on a virtual target. In a separate lab workflow, I built a Python pipeline for extracting structural vectors from packaging artwork, supporting generation of more than 20,000 training inputs. At the Koita Centre for Digital Health, I also contributed to a wearable-sensor pipeline achieving approximately 70\% sensitivity for asthma exacerbation detection in study evaluation, building preprocessing and change-point features with patient-wise evaluation.\par\vspace{10pt}
I would bring hands-on Python and C++ experience, an interest in model behaviour under imperfect inputs, and a willingness to work through the details of calibration and evaluation. The opportunity to develop AI systems alongside experienced Honeywell engineers would let me build on this work in real-world automation projects. Thank you for considering my application.\par\vspace{14pt}
Yours sincerely,\\[5pt]
Aarush Ram Anandh
\end{document}
